\begin{ouverture}[trigo]{2}{Module, arguments et plan complexe}{Nombres complexes}{Affixes · Distances · Cercle unité · Forme trigonométrique}
\objectif{Passer d’un point à son affixe et inversement.}
\objectif{Calculer modules et arguments.}
\objectif{Utiliser les propriétés multiplicatives du module.}
\objectif{Reconnaître des ensembles de points simples.}
\prerequis{Coordonnées, distance euclidienne et cercle trigonométrique.}
\end{ouverture}
\essentiel

\essSection{1}{L'affixe encode un point ou un vecteur}
Dans un repère orthonormé direct, le point $M(x,y)$ a pour affixe $z=x+\mathrm iy$.
Si $A,B$ ont pour affixes $a,b$, le vecteur $\vect{AB}$ a pour affixe $b-a$ et le milieu de $[AB]$ a pour affixe $(a+b)/2$.
Le conjugué correspond à la symétrie par rapport à l'axe réel ; l'opposé à la symétrie de centre $O$.
\begin{definition}[Module]
$|z|=\sqrt{x^2+y^2}$ est la distance $OM$.
En particulier $AB=|b-a|$ et $|z|=0$ si et seulement si $z=0$.
\end{definition}
\figurelivre{e02-module.pdf}{Pour $z=3+2\mathrm i$, le module est la longueur $OM=\sqrt{13}$. Un argument est l'angle orienté du demi-axe réel positif vers $\vect{OM}$.}
\essSection{2}{Les identités du module}
\begin{propriete}
$|z|^2=z\bar z$, $|zw|=|z||w|$, $|z^n|=|z|^n$.
Pour $z\ne0$, $|1/z|=1/|z|$ ; pour $w\ne0$, $|z/w|=|z|/|w|$.
\end{propriete}
\begin{demonstration}
$z\bar z=(x+\mathrm iy)(x-\mathrm iy)=x^2+y^2=|z|^2$.
Donc $|zw|^2=zw\bar z\bar w=(z\bar z)(w\bar w)=|z|^2|w|^2$. Les modules étant positifs ou nuls, prendre la racine donne la règle du produit.
Une récurrence donne la puissance ; le produit $z(1/z)=1$ donne la règle de l'inverse.
\end{demonstration}
\essSection{3}{L'argument est défini modulo $2\pi$}
\begin{definition}
Si $z\ne0$, un argument $\theta$ vérifie $z=|z|(\cos\theta+\mathrm i\sin\theta)$.
Tous les arguments sont $\theta+2k\pi$, $k\in\Z$.
Le nombre $0$ n'a pas d'argument.
\end{definition}
Pour identifier un angle, utiliser simultanément $\cos\theta=x/|z|$ et $\sin\theta=y/|z|$. La tangente seule perd l'information sur le quadrant.
\begin{exemple}
$z=-1+\mathrm i\sqrt3$ a pour module $2$, avec cosinus $-1/2$ et sinus $\sqrt3/2$ : un argument est $2\pi/3$.
Ainsi $z=2(\cos(2\pi/3)+\mathrm i\sin(2\pi/3))$.
\end{exemple}
\essSection{4}{Le cercle unité et les lieux}
L'ensemble $\mathbb U=\{z\in\C:|z|=1\}$ est stable par produit et par inverse, car les modules restent $1$.
Sur $\mathbb U$, $1/z=\bar z$.
$|z-a|=R$ décrit un cercle de centre d'affixe $a$ et de rayon $R>0$ ; $|z-a|\leq R$ décrit le disque fermé.
L'égalité $|z-a|=|z-b|$, pour $a\ne b$, décrit la médiatrice de $[AB]$.
\figurelivre{e02-lieux.pdf}{À gauche : $|z-(1+\mathrm i)|=2$. À droite : $|z+1|=|z-1|$, médiatrice de $[-1,1]$ sur l'axe réel.}
\essSection{5}{Une inégalité géométrique, en approfondissement}
L'inégalité triangulaire $|z+w|\leq|z|+|w|$ traduit que la distance directe ne dépasse pas la longueur d'un trajet brisé.
Algébriquement, $|z+w|^2=|z|^2+|w|^2+2\Re(z\bar w)$ et $\Re(z\bar w)\leq|z\bar w|=|z||w|$.
L'égalité correspond à des vecteurs de même direction et de même sens, avec les cas où l'un est nul.

\exercices
\markright{Prendre confiance}
\begingroup\renewcommand{\otBesoinPlace}[1]{}\begin{minipage}{\linewidth}
\rubrique{Prendre confiance}
\begin{exercice}[1]{Affixes}

\begin{questions}
\item Donner l’affixe de $M(-2,3)$.
\item Donner les coordonnées du point d’affixe $4-\mathrm i$.
\item Si $a=1+\mathrm i$, $b=3+4\mathrm i$, donner l’affixe de $\vect{AB}$.
\end{questions}
\end{exercice}
\begin{exercice}[1]{Modules}

\begin{questions}
\item Calculer $|3-4\mathrm i|$.
\item Calculer $|2\mathrm i|$ et $|-5|$.
\item Calculer $|(1+\mathrm i)^5|$.
\end{questions}
\end{exercice}
\begin{exercice}[1]{Arguments remarquables}

\begin{questions}
\item Donner les arguments de $1$, puis de $\mathrm i$.
\item Donner ceux de $-1$.
\item Pourquoi zéro n’en a-t-il pas ?
\end{questions}
\end{exercice}
\begin{exercice}[1]{Forme trigonométrique}

\begin{questions}
\item Mettre $1+\mathrm i$ sous forme trigonométrique.
\item Faire de même pour $-\sqrt3+\mathrm i$.
\item Revenir à la forme algébrique de $2(\cos(\pi/6)+\mathrm i\sin(\pi/6))$.
\end{questions}
\end{exercice}
\end{minipage}\endgroup
\clearpage
\markright{Appliquer et relier}
\begingroup\renewcommand{\otBesoinPlace}[1]{}\begin{minipage}{\linewidth}
\rubrique{Appliquer et relier}
\begin{exercice}[2]{Distances et milieu}
$a=1+2\mathrm i$ et $b=5-\mathrm i$.
\begin{questions}
\item Calculer $AB$.
\item Donner l’affixe du milieu.
\item Donner l’affixe du symétrique de $A$ par rapport à ce milieu.
\end{questions}
\end{exercice}
\begin{exercice}[2]{Le cercle unité}
$z=(3+4\mathrm i)/5$.
\begin{questions}
\item Vérifier $z\in\mathbb U$.
\item Calculer $1/z$.
\item Justifier que $z^{10}\in\mathbb U$ sans développer.
\end{questions}
\end{exercice}
\begin{exercice}[2]{Décrire des lieux}

\begin{questions}
\item Décrire $|z-2+\mathrm i|=3$.
\item Décrire $|z|\leq2$.
\item Décrire $|z-1|=|z+1|$.
\end{questions}
\end{exercice}
\begin{exercice}[2]{Un argument mal choisi}
Un élève dit que $-1-\mathrm i$ a pour argument $\pi/4$ car le quotient des coordonnées vaut $1$.
\begin{questions}
\item Calculer le module.
\item Donner un argument correct.
\item Expliquer l’erreur.
\end{questions}
\end{exercice}
\end{minipage}\endgroup
\clearpage
\markright{Chercher et argumenter}
\begingroup\renewcommand{\otBesoinPlace}[1]{}\begin{minipage}{\linewidth}
\rubrique{Chercher et argumenter}
\begin{exercice}[3]{Démontrer le produit}

\begin{questions}
\item Écrire $|zw|^2$ à l’aide des conjugués.
\item Regrouper les facteurs.
\item Pourquoi peut-on prendre la racine sans ambiguïté de signe ?
\end{questions}
\end{exercice}
\begin{exercice}[3]{Médiatrice par le calcul}
$a=-1$, $b=3$ et $z=x+\mathrm iy$.
\begin{questions}
\item Écrire l’égalité $|z-a|=|z-b|$ avec $x,y$.
\item La simplifier.
\item Relier la réponse à la géométrie.
\end{questions}
\end{exercice}
\begin{exercice}[3]{Triangle et égalité}
$z=1+\mathrm i$ et $w=2+2\mathrm i$.
\begin{questions}
\item Comparer $|z+w|$ et $|z|+|w|$.
\item Expliquer géométriquement l’égalité.
\item Que se passe-t-il si l’on remplace $w$ par $-z$ ?
\end{questions}
\end{exercice}
\begin{exercice}[3]{Module constant ?}
$z=(x-\mathrm i)/(x+\mathrm i)$ avec $x\in\R$.
\begin{questions}
\item Vérifier que le dénominateur est non nul.
\item Prouver $|z|=1$.
\item Le point d’affixe $1$ peut-il être obtenu pour un réel $x$ ?
\end{questions}
\end{exercice}
\end{minipage}\endgroup
